\begin{table}[t]
\centering
\scriptsize
\setlength{\tabcolsep}{3.5pt}
\resizebox{\linewidth}{!}{\begin{tabular}{lrrrrrr}
\toprule
Model & NLL $\downarrow$ & cell NLL $\downarrow$ & MMD $\downarrow$ & AUC $\downarrow$ & $r_\mu$ $\uparrow$ & Type acc. $\uparrow$ \\
\midrule
Global mean & 6.732 (+1.125 [+0.939, +1.326]) & 6.732 & 0.8329 & 1.000 & 0.553 & 0.038 \\
Cell-type mean & \textbf{5.606} & \textbf{5.606} & 0.1816 & 0.985 & \textbf{0.931} & \textbf{0.879} \\
\rowcolor{ours}CAD (ours) & 5.637 (+0.031 [+0.001, +0.076]) & 5.637 & \textbf{0.1732} & \textbf{0.973} & 0.927 & 0.791 \\
\bottomrule
\end{tabular}}
\caption{Human PBMC (Hao et al.), test split: 24 held-out cell types. NLL: multinomial negative log-likelihood per UMI of the pooled real cells under the model (Delta: paired difference to Cell-type mean [95\% bootstrap over cell types]). MMD, AUC (real vs generated classifier; 0.5 is ideal), $r_\mu$ (gene-mean correlation), type acc.: a classifier trained on real cells recovers the requested type use cells from the shared Dirichlet-multinomial sampler. Best in bold; our rows shaded.}
\label{tab:pbmc_test_cond}
\end{table}